\input{./._relpath-to-latexroot.ltx}
\documentclass[\latexroot/\projectname]{subfiles}
\input{\latexroot/@resources/subfile-setup.ltx}
\begin{document}

\section{Conclusions}\whenintegrated{\label{conclusion}}
\whenintegrated{\label{sec:conclusion}}

We solved the fixed-price household problem and checked the saving policies at the upper end of the asset grid.

Workers solve a standard consumption-saving problem, while entrepreneurs also choose business capital subject to an enforcement constraint and a discrete capital-grid ceiling. Occupational choice compares the value of these two branches. In the baseline run, the young-only entrepreneur share is 8.33 percent, close to the original paper's 7.5 percent benchmark. The paper's benchmark also includes older households.

The notebook's young-only wealth Gini is 0.67, compared with 0.80 in the paper's baseline model, and the top 1 percent wealth share is 14.4 percent rather than 31 percent. These notebook figures are PE diagnostics, not formal bounds on the paper's Table~6 moments. The young-only distribution omits aging and inheritance flows. The household problem includes descendant values, and prices stay fixed throughout the calculation.

The most natural next step would be to replace the young-only stationary distribution check with a full demographic distribution over young workers, young entrepreneurs, old entrepreneurs, and retirees. After that, the model could be closed with market clearing and social-security budget balance, and only then would it be appropriate to compare the reproduced wealth distribution and counterfactuals more directly with the original paper's quantitative tables.

\smartbib

\end{document}
